%% kapa: exam=1 round=22 year=2011 period=1 form=A subject=economics subject_ko=경제원론 seq=25 number=65
%% source: Q-Net 감정평가사 기출문제 (cst00309) artlSeq=5036790 file=extracted/1008/5036790/2011년 제22회 감정평가사 1차시험 기출문제/1교시.hwp
%% conversion: kordoc 4.24.51 -> tools/gen1.py (auto); review=proofread-sonnet; flags=none
\begin{question}{65}
IS곡선에 관련된 설명으로 옳지 않은 것은? (단, IS곡선은 우하향)
\begin{choices}
  \item IS곡선은 생산물시장의 균형을 이루는 이자율과 국민소득의 조합을 나타낸다.
  \item 현재의 이자율과 국민소득의 조합점이 IS곡선보다 위쪽에 있다면, 생산물시장에서 수요가 공급을 초과하고 있음을 의미한다.
  \item 조세부담이 증가하면 IS곡선은 좌측으로 이동한다.
  \item 정부의 재정지출이 증가하면 IS곡선은 우측으로 이동한다.
  \item 한계소비성향이 높아질수록 IS곡선은 더 완만해진다.
\end{choices}
\end{question}
\answer{2}
